\documentclass[11pt]{article}

\usepackage[utf8]{inputenc}
\usepackage[T1]{fontenc}
\usepackage{lmodern}
\usepackage{geometry}
\usepackage{amsmath,amssymb,amsfonts,amsthm,mathtools}
\usepackage{bm}
\usepackage{enumitem}
\usepackage{microtype}
\usepackage{hyperref}
\usepackage{titlesec}
\usepackage{booktabs}
\usepackage{array}
\usepackage{setspace}
\usepackage{mathrsfs}
\usepackage{physics}

\geometry{margin=1in}
\hypersetup{colorlinks=true, linkcolor=black, urlcolor=blue, citecolor=black}
\setlist[enumerate]{topsep=0.4em,itemsep=0.35em}
\setlist[itemize]{topsep=0.3em,itemsep=0.25em}
\titleformat{\section}{\large\bfseries}{\thesection.}{0.5em}{}
\titleformat{\subsection}{\normalsize\bfseries}{\thesubsection.}{0.5em}{}
\setstretch{1.06}

\newcommand{\Aether}{\AE{}ther}
\newcommand{\Aflow}{\AE{}ther-flow}
\newcommand{\TheoryName}{\Aflow{} Interpretation of Relativity}
\newcommand{\TheoryProgram}{\Aether{} / \Aflow{} framework}
\newcommand{\CompletedMetric}{g^{\sharp}}
\newcommand{\NanoSectionPackage}{\mathfrak X^{\mathrm{nanosec}}}
\newcommand{\PicoSectionPackage}{\mathfrak Y^{\mathrm{picosec}}}
\newcommand{\PicoMinSectionCore}{\mathfrak Y^{\mathrm{pic,sec}}}
\newcommand{\FemtoSectionPackage}{\mathfrak Z^{\mathrm{femtosec}}}
\newcommand{\AttoSectionPackage}{\mathfrak A^{\mathrm{attosec}}}

\newtheorem{definition}{Definition}
\newtheorem{proposition}{Proposition}
\newtheorem{remark}{Remark}
\newtheorem{corollary}{Corollary}
\newtheorem{theorem}{Theorem}

\title{\textbf{The \TheoryName{}}\\[0.4em]
\large Primitive-Reservoir Observer-Reduction\\
\large Femto-Section-Generating Substrate Atto-Section Dynamics\\
\large Next Benchmark Criterion}
\author{}
\date{}

\begin{document}

\maketitle

\begin{abstract}
The current active-primary theorem already removed the current femto-section package \(\FemtoSectionPackage\) as a primitive stopping point by deriving it canonically from the deeper current atto-section package \(\AttoSectionPackage\). This renews the live frontier. Another same-output deeper-origin relay that merely preserves \(\AttoSectionPackage\) and therefore merely reproduces the same current femto-section package \(\FemtoSectionPackage\), the same current pico-section package \(\PicoSectionPackage\), the same current minimizer-section core \(\PicoMinSectionCore\), the same current nano-section package \(\NanoSectionPackage\), and the same downstream micro-section and sub-section packages no longer counts as the honest default next benchmark-facing manuscript. Neither does a re-expansion back to the larger minimizer-image-core or full pico-section presentation, and neither does a quantum branch that does not do that same benchmark-facing work.

This note records the routing consequence of that theorem. The next honest benchmark-facing move must either derive or justify the current femto-section-generating substrate atto-section dynamics package \(\AttoSectionPackage\) from still more primitive substrate structure in a way that changes benchmark-facing status, or prove a genuinely smaller theorem-level collapse strictly inside \(\AttoSectionPackage\). The benchmark boundary remains fixed throughout. The one operative metric is still exactly \(\CompletedMetric\), the ordinary matter route is unchanged, there is no second operative metric, the low-energy matter law is not enlarged, and no stronger promotion language is licensed. The positive result is therefore routing discipline at the current atto-section-dynamics frontier.
\end{abstract}

\section{Introduction}

The active derivational program of \TheoryName{} is no longer at the current femto-section package as its default stopping point. The current active-primary theorem already showed that the current femto-section package \(\FemtoSectionPackage\) is canonically generated by the deeper current atto-section package \(\AttoSectionPackage\). That theorem therefore spent the package-derivation option that had remained open below the current femto-section frontier.

The present note records the routing consequence of that gain. It does not reopen the already-generated current femto-section package, and it does not claim a completed first-principles substrate derivation of gravity. It records only which kind of next manuscript would now count as an honest benchmark-facing gain once the live burden has moved back to the deeper current atto-section package itself.

\paragraph{Framework and claim boundary.}
\TheoryProgram{} denotes the broader \Aether{} / \Aflow{} framework built on the claims that the \Aether{} is the underlying four-dimensional substrate of reality, the \Aflow{} is its intrinsic ordered motion, observed three-dimensional space is the local experiential slice of that deeper substrate, S-time is the experienced order of change arising from matter, light, and the \Aflow{}, and observed expansion is the three-dimensional appearance of deeper four-dimensional ordered motion. Within that framework, \TheoryName{} denotes the exact-closure theory in which the effective gravitational dynamics are adopted to be exactly Einsteinian with universal matter coupling. In the active sequence, \emph{adoption} means use of that established relativistic dynamics without claiming substrate derivation, while \emph{derivation} is reserved for a first-principles recovery from explicit substrate variables.

The present note stays strictly inside that boundary. It records only which kind of next manuscript would now count as an honest benchmark-facing gain at the current atto-section-dynamics frontier.

\section{Fixed Inputs}

\begin{definition}[Atto-section frontier inputs]
The criterion recorded here uses the following fixed inputs.
\begin{enumerate}
    \item The benchmark exact-closure package, which fixes the public one-metric GR target of the repository.
    \item The benchmark gatekeeping layer, including the recorded safeguard that blocks same-output deeper-origin relays from counting as new front-line progress once the live benchmark-relevant datum is unchanged.
    \item The prior renewed pico-section-package frontier and the theorem deriving the current pico-section package from deeper pico-section-generating substrate femto-section dynamics, which had already moved the live burden below the current pico-section package.
    \item The current active-primary theorem, which proves that the current femto-section package \(\FemtoSectionPackage\) is canonically generated by the deeper current femto-section-generating substrate atto-section dynamics package \(\AttoSectionPackage\).
\end{enumerate}
\end{definition}

\begin{remark}
The present note does not replace those manuscripts. It records the routing consequence of reading them together under the active safeguard against recursive depth drift.
\end{remark}

\section{Why the Current Burden Returns to the Atto-Section Package}

\begin{proposition}[The live burden is renewed at atto-section-package scope]
At the current stage of the primitive-reservoir observer-reduction line, the current active-primary theorem renews the live benchmark-facing burden at the level of the current atto-section package \(\AttoSectionPackage\).
\end{proposition}

\begin{proof}
That theorem changes benchmark-facing status because it removes the current femto-section package \(\FemtoSectionPackage\) as the primitive stopping point. Once \(\FemtoSectionPackage\) is shown to be canonically induced by \(\AttoSectionPackage\), the live burden no longer sits on the current femto-section package itself. It returns to the deeper package from which that package is generated. The earlier pico-section-package theorem and the preserved renewed pico-section-package routing handoff remain main-line history and routing support, but they no longer define the live frontier once the current femto-section package is itself derived from \(\AttoSectionPackage\).
\end{proof}

\begin{definition}[Benchmark-neutral deeper relay at current atto-section scope]
A \emph{benchmark-neutral deeper relay at current atto-section scope} is a manuscript that preserves the same benchmark one-metric output, the same ordinary matter route, the same current atto-section package \(\AttoSectionPackage\), the same current femto-section package \(\FemtoSectionPackage\), the same current pico-section package \(\PicoSectionPackage\), the same current minimizer-section core \(\PicoMinSectionCore\), the same current nano-section package \(\NanoSectionPackage\), and the same downstream micro-section and sub-section packages while merely relocating the unexplained substrate layer deeper, repackaging the same current atto-section package, or re-expanding the presentation back to the larger minimizer-image-core or full pico-section package without a new compression, necessity, uniqueness, or comparably sharp routing gain.
\end{definition}

\section{The Current Atto-Section Next Benchmark Criterion}

\begin{definition}[Admissible next benchmark-facing manuscript at current atto-section scope]
At the current atto-section-dynamics frontier, a manuscript counts as an admissible next benchmark-facing move only if it does at least one of the following:
\begin{enumerate}
    \item derives or justifies the current femto-section-generating substrate atto-section dynamics package \(\AttoSectionPackage\) from still more primitive substrate structure in a way that does more than merely relocate the unexplained layer deeper;
    \item proves a genuinely smaller theorem-level collapse strictly inside \(\AttoSectionPackage\);
    \item does either of the previous two while preserving the same benchmark one-metric package, the same ordinary matter route, and the same claim boundary.
\end{enumerate}
\end{definition}

\begin{theorem}[Next honest benchmark-facing task at current atto-section scope]
The next honest benchmark-facing manuscript below the current atto-section-dynamics frontier must either derive or justify the current femto-section-generating substrate atto-section dynamics \(\AttoSectionPackage\) from still more primitive substrate structure in a benchmark-relevant way, or else prove a genuinely smaller theorem-level collapse strictly inside that deeper package. A manuscript that only repeats the same current atto-section package through another deeper relay, re-expands the discussion back to the larger minimizer-image-core or full pico-section presentation without such a new gain, or invokes a quantum branch without doing the same benchmark-facing work does not satisfy the criterion.
\end{theorem}

\begin{proof}
By Proposition 1, the live burden already lies at the deeper current atto-section package \(\AttoSectionPackage\) rather than at the current femto-section package \(\FemtoSectionPackage\) or the earlier pico-section-package handoff. Therefore a subsequent manuscript must improve on that status rather than merely accompany it. A deeper-origin manuscript can do so only if it changes benchmark-facing status by more than depth alone. If it simply reproduces the same current atto-section package and its same downstream outputs, the routing safeguard classifies it as benchmark neutral. The remaining honest alternatives are therefore exactly those stated in the definition: either a more primitive derivation or justification of the current atto-section package itself, or a sharper internal collapse strictly inside that package.
\end{proof}

\section{What the Next Move Should Not Be}

\begin{proposition}[Screened next moves at the current atto-section-dynamics frontier]
At the current atto-section-dynamics frontier, none of the following is the default next benchmark-facing move:
\begin{enumerate}
    \item another same-output deeper-origin relay that merely preserves the current atto-section package \(\AttoSectionPackage\) and therefore merely reproduces the same current femto-section package \(\FemtoSectionPackage\), the same current pico-section package \(\PicoSectionPackage\), the same current minimizer-section core \(\PicoMinSectionCore\), the same current nano-section package \(\NanoSectionPackage\), and the same downstream micro-section and sub-section packages;
    \item another re-expansion back to the larger minimizer-image-core or the full pico-section presentation as if those larger presentations were again independently live burdens;
    \item a quantum branch that does not derive or justify the same deeper atto-section package, collapse it further, or close a benchmark derivation gate in a genuinely new way;
    \item a manuscript that introduces a second operative metric, enlarges the ordinary matter law, or uses stronger promotion language.
\end{enumerate}
\end{proposition}

\begin{proof}
Item (1) fails the preceding criterion because it leaves the renewed live burden unchanged. Item (2) likewise fails because it re-expands to larger already-compressed presentations rather than removing remaining freedom in \(\AttoSectionPackage\) or proving a smaller internal datum there. Item (3) does not directly address the live burden at current scope unless it performs the same benchmark-facing work named above. Item (4) violates the fixed claim boundary of the benchmark package and therefore cannot count as a disciplined next step on the active line. The benchmark package remains the exact one-metric GR package already fixed by adoption, so any admissible next manuscript must preserve that boundary.
\end{proof}

\begin{remark}
This screening statement is not a denial that those deeper or alternative manuscripts may matter strategically. It is a routing statement: they are not the default way to spend the next benchmark-facing move from the current atto-section-dynamics frontier unless they do the same benchmark-facing work named above.
\end{remark}

\section{Conclusion}

The positive result of this note is routing discipline at the current atto-section-dynamics frontier. The current active-primary theorem remains the live benchmark-facing gain because it already removes the current femto-section package as the unexplained stopping point on the active line. The earlier pico-section-package theorem and renewed pico-section-package routing handoff remain preserved main-line history, but under the recursive-depth safeguard they are no longer themselves the default current frontier.

The scope remains narrow and honest. The benchmark package is unchanged: the one operative metric is still exactly \(\CompletedMetric\), the ordinary matter route is unchanged, there is no second operative metric, and the low-energy matter law is not enlarged. Nothing here upgrades adoption to derivation or licenses stronger promotion language.

The next open burden is therefore clear. The next benchmark-facing manuscript should derive or justify the current femto-section-generating substrate atto-section dynamics from still more primitive substrate structure in a way that changes benchmark-facing status, unless the sharper honest gain available instead is a genuinely smaller theorem-level collapse strictly inside that deeper package. Another same-output deeper relay that merely preserves the current atto-section package, a re-expansion back to the larger minimizer-image-core or full pico-section presentation, or a quantum branch that does not do that same benchmark-facing work is not the clean next manuscript target at current scope.

\end{document}
